\chapter{Reprodutibilidade e acompanhamento}
Esta versão é a proposta de pesquisa, e não uma dissertação concluída.
O repositório público \url{https://github.com/flavioluiz/TGdoWayne} reúne o TG
original, o template fornecido, as fontes deste documento, o plano de execução
e os resultados das checagens preliminares.

O arquivo \texttt{README.md} registra a etapa atual, as entregas concluídas,
a próxima etapa e o endereço do PDF da versão mais recente. O plano geral
está em \path{implementation_plan/README.md}; cada arquivo na subpasta
\texttt{commits} detalha um marco. A sequência prevê 13 commits e 13 versões:
proposta; introdução; revisão; fundamentos; resposta de PTA; simulações;
inferência; compressão; robustez; helicidade zero; aplicação ou extensão
simulada; discussão e artigo; auditoria final.

Cada commit de marco deverá conter as fontes, o PDF cumulativo, as notas da
versão e um manifesto de integridade. A tag identifica o mesmo commit que
contém o PDF; a publicação no GitHub anexa esse arquivo sem recompilá-lo.
As versões a partir de C02 constituirão a dissertação em desenvolvimento.
Capítulos não concluídos serão identificados como planejados, sem resultados
fictícios. Correções após publicação receberão nova versão e novo PDF.

A etapa C01 contém um script em Python que usa apenas a biblioteca padrão.
A partir da raiz do repositório, os comandos são:
\begin{quote}\small\ttfamily
python3 scripts/checagens\_preliminares.py\par
python3 scripts/build\_release.py --build-only\par
python3 scripts/verify\_release.py v0.1.0
\end{quote}
O segundo comando requer pdfLaTeX, latexmk e Biber, além dos pacotes descritos
no README. O arquivo JSON com as checagens contém aritmética racional para os
vínculos e os testes de curvatura, e aritmética decimal de alta precisão para
as escalas físicas. Nenhum dado observacional foi ajustado nesta versão.

